\subsection{Full Factorial and Search-Retention Evidence}
\label{app:factorial-details}

The 49-query control holds 277 slots and its template inventory fixed within decoder comparisons, distinct from the 331-query matched study. The known excluded query shares a delexicalized text group with earlier training/development data; grouping uses transitive five-gram Jaccard similarity at least 0.80. Exclusion addresses that known crossing, not prior benchmark exposure or unrecognized semantic duplicates.

Main Section~\ref{M-subsec:method_controls} specifies the 24-cell design; Table~\ref{M-tab:grounding-factorial-diagnostic} reports every one-to-one control. All 24 cells are preserved in the release's numeric ledgers and score/choice projections. Candidate identities, source context, values and lexical traces remain in local originals, not the public projection. The ordinal control's cardinalization and partially assisted inventory limit interpretation.

The historical contextual decoder intervals and unchanged $5/49$ full-tolerance count are reported in main Section~\ref{M-subsec:grounding-followup}. The separate 240-test family covers every one-factor contrast and four endpoints; some contextual-versus-ordinal slot contrasts survive Holm correction, but no whole-instance endpoint does. No interval is promoted into a confirmatory superiority or equivalence claim.

\medskip\noindent\textbf{Candidate availability, pruning and objective quality.}
The earlier 50-query joint-grounding pilot selects 3 tolerance-correct
instances with either local or joint scoring. In the saved common pools,
32/50 permit a criterion-correct complete mapping, but only 8/50 retain one
in the top ten. Reconciliation with identical saved beam outputs separates
14 cases losing correct prefixes, 6 generating then pruning a complete
correct mapping, 4 losing it at the final top-ten cut, 5 retaining it below
the selected mapping, and 3 selecting it. The other 18 cases lack a complete
compatible mapping in the pool or inventory. Width and ranking limitations
must therefore be distinguished from candidate absence. Optimality claims
about beam methods also depend on their search restriction
\cite{huang2017beamfollowup}; that neural-generation result is methodological
context, not a guarantee for this grounding algorithm.

An exact MILP diagnostic holds the cached local/pairwise objective and
domains fixed. Among the 32 pool-compatible cases, 23 have a strictly
higher-scoring incorrect optimum than their best criterion-correct mapping;
6 have tied correct and incorrect optima, and 3 select a correct optimum.
The correctness-restricted solves are oracle diagnostics, never inference
procedures. Exact mention-based top-ten enumeration retains a correct
mapping in 9/50 cases versus 8/50 for the beam, with tied-optimum membership
qualifications. More search cannot repair the 23 strict objective failures.
Ten of the 50 exact top-ten lists contain duplicate numerical/type mappings
realized by different mention indices. A separate 49-query follow-up ranks without evaluation values, conditional on the same supplied inventory. It reuses
40 already-distinct lists and performs 23 additional solves for the other
nine. Every retained top-ten list is certified, but criterion-correct list
retention remains 9/49, with zero gained or lost queries; selected correctness
remains 3/49. Thus removing these duplicate mappings yields no observed
correctness gain for this frozen objective. Mention equality, numerical equality and semantic identity are
different concepts. No gold value is used to rank inference candidates.

Detailed source/reference conventions, the versioned correction overlay and the historical tolerance limitation are given in Section~\ref{app:reference-conventions}. No reference correction is used during inference, and no new hard unit or mathematical relation was validated by these studies.
